\documentclass[10pt,a4paper]{amsart}
\usepackage[margin=19mm]{geometry}
\usepackage{amsmath,amssymb,mathtools}
\usepackage[scr=rsfs,cal=euler]{mathalfa}
\usepackage{xfrac}
\usepackage[unicode,hidelinks,bookmarks=false]{hyperref}
\newcommand{\ie}{i.e.\ }
\newcommand{\cf}{cf.\ }
\newcommand{\resp}{respectively\ }
\newcommand{\Subsection}{\subsection}
\providecommand{\nobreakdash}{\nobreak\hbox{-}}
\setlength{\emergencystretch}{2em}
\multlinegap0pt
\usepackage{bm}
\usepackage{bbm}
\usepackage{dsfont}
\usepackage{xspace}
\usepackage{cancel}
\usepackage{mathtools}
\usepackage{amsthm}
\makeatletter
\@ifundefined{definition}{\newtheorem{definition}{Definition}}{}
\@ifundefined{theorem}{\newtheorem{theorem}{Theorem}}{}
\makeatother
\title[Marcello's completion of graphs]{Marcello's completion of graphs}
\author{Johan Kok}
\date{Source: arXiv:2507.02015v1 (CC BY 4.0)}
\begin{document}
\maketitle

\noindent\emph{Source and licence.} Marcello's completion of graphs. Version arXiv:2507.02015v1 (CC BY 4.0), distributed under the Creative Commons Attribution 4.0 International licence, as recorded in the arXiv licence metadata for this exact version and shown on the abstract page https://arxiv.org/abs/2507.02015v1. Full rights notice: licenses/arxiv-2507.02015-CC-BY-4.0.txt.\par\smallskip

\noindent\textit{Editorial scope.} Counted unit: the Theorem whose source label is \texttt{thm2} in the article (source0.tex lines 268--270, source label \texttt{thm2}), delivered together with the article's own complete original proof of it (source0.tex lines 271--307, one \texttt{proof} environment of 37 lines). No mathematics is written, repaired or re-derived: statement and proof are the article's own text line for line. Required context retained uncounted: def1 (lines 58--68). Every definition, display and label of those lines is retained unchanged. Omitted from this package and not redistributed: 1--57, 69--267, 308--559. Nothing inside a retained excerpt is omitted or altered: every retained body is delivered exactly as the article's lines are written, so no omission receipt of any kind accompanies this package. Citations: the retained excerpts carry no citation command at all, so no bibliography entry of the article is transcribed and no reference is redistributed. Reference adaptation: none was needed and none was made; every reference inside the delivered unit resolves inside it, so no unresolved reference or citation is produced. Printed slips kept literal: no printed word, symbol, hypothesis or number of the counted statement or of its proof was altered. Layout and class: the article's own class is \texttt{amsart} with packages including amsmath, amssymb, amsfonts, amsthm, graphicx, enumerate, tikz, eucal, float; the wrapper is the lane's pinned \texttt{amsart} editorial wrapper at 10pt on A4, which restates the theorem-like environments the retained text needs and adds the article's own macro definitions that the retained text needs (none), with extra wrapper packages (bm, bbm, dsfont, xspace, cancel, mathtools). The delivered numbering is the unit's own. Assets: no retained span contains a figure environment, a graphics-inclusion command, a PDF-page inclusion, a table or a TikZ picture; no image file is copied into this package, the package declares no asset dependency and no image file is redistributed with it. Licence: version arXiv:2507.02015v1 of this article is distributed under the Creative Commons Attribution 4.0 International licence, as recorded in the arXiv licence metadata for this exact version and shown on the abstract page https://arxiv.org/abs/2507.02015v1. A full rights notice is delivered at licenses/arxiv-2507.02015-CC-BY-4.0.txt. Characters: every retained excerpt is pure ASCII, so the delivered engine, XeLaTeX with its default Latin Modern text font, reports no missing character.
\begin{definition}\label{def1}
Let $H_{i-1} = G_{i-1}$. Let $X_{i-1} \subseteq V(G_{i-1})$ where $X_{i-1} = \{v_j: 0 < deg_{G_{i-1}}(v_j) < n-1\}$ and $\ell = |X_{i-1}|$. 
 
For $k = 1,2,3,\dots, \ell$ go to Step 1: 

Step 1:  Let $j = k$ and select any vertex $v_{t_j} \in X_{i-1}$, ($t_j = s$ for some $v_s \in X_{i-1}$). Go to Step 2.

Step 2: If possible, add at most $deg_{G_{i-1}}(v_{t_j})$ edges (called, Marcello edges) to non-neighbors of $v_{t_j}$ in $H_{j-1}$. Label the resultant graph $H_j$. If $H_j$ is complete then label it as $G_i$ and exit ($G$ has been completed). Else, $X_{i-1}\mapsto X_{i-1}\backslash \{v_{t_j}\}$. Go to step 3.

Step 3: If $X_{i-1} = \emptyset$ or if all vertices remaining, say $v_q \in X_{i-1}$ has $deg_{H_j}(v_q) = n-1$ then label it as $G_i$ and exit (a global iteration completed). Else, consider the next $k$ and go to Step 1 (a local iteration completed).\\
\end{definition}

\begin{theorem}\label{thm2}
Let $G = K_n\cup K_m$, $n \geq 2$, $m \geq 1$. Then $\varpi(G) = 1$.
\end{theorem}

\begin{proof}
Let $V(K_n) = \{v_i: i = 1,2,3,\dots,n\}$ and $V(K_m) = \{u_j:j = 1,2,3,\dots,m\}$. Without loss of generality assume $n \geq m$.

Case 1: Let $n > m$. For this case $n-m \geq 1$. By Definition \ref{def1} the following edges can be added on the first global iteration:
\begin{center}
$v^+_1u_j$ for $j = 1,2,3,\dots,m$;

$v^+_2u_j$ for $j = 1,2,3,\dots,m$;

$v^+_3u_j$ for $j = 1,2,3,\dots,m$;

$\vdots$

$v^+_nu_j$ for $j = 1,2,3,\dots,m$.
\end{center}
Clearly, $G_1$ is complete. This settles the result for $n > m$.

Case 2(a): Let $n = m$ and even. By Definition \ref{def1} the edges which can be added on the first global iteration to exhaust the vertices $v_i$, $ i = 1,2,3,\dots,n$ are: 
\begin{center}
$v^+_1u_j$ for $j = 1,2,3,\dots,m-1$;

$v^+_2u_j$ for $j = 2,3,4,\dots,m$;

$v^+_3u_j$ for $j = 1,2,3,\dots,m-1$;

$v^+_4u_j$ for $j = 2,3,4,\dots,m$;

$\vdots$

$v^+_{n-1}u_j$ for $j = 1,2,3,\dots,m-1$;

$v^+_nu_j$ for $j = 2,3,4,\dots,m$.
\end{center}
Clearly, the vertices $u_1,u_m \in V(K_m)$ have not been \lq completed\rq. Each of these vertices may initiate $m-1$ Marcello edges and exactly $\frac{m}{2}$ Marcello edges are required. We determine for which values of $m$ we have $m-1 \geq \frac{m}{2}$. The latter is valid for $m \geq 2$. Hence, Marcello completion is possible in one global iteration. Therefore, $\varpi(G) = 1$.

Case 2(b): For $n = m$ and odd, the proof follows by similar reasoning as that found in the proof of Case 2(a).
\end{proof}

\end{document}
